problem:  
Let \((P, \mathcal{H}, g)\) be a compact contact sub‑Riemannian manifold, where \(\mathcal{H} = \ker \alpha\) for a contact 1‑form \(\alpha\) and \(g\) is a smooth metric on \(\mathcal{H}\).  
Suppose a compact connected Lie group \(G\) acts smoothly and properly on \(P\), preserving \(\alpha\) and \(g\), with infinitesimal action contained in \(\mathcal{H}\). Denote by \(\pi : P \to M := P/G\) the quotient, which is assumed to be a smooth manifold, and by \(\mathcal{H}_{\mathrm{red}}\) the push‑forward of \(\mathcal{H}\) under \(\pi\). The quotient inherits a natural sub‑Riemannian structure \((M, \mathcal{H}_{\mathrm{red}}, g_{\mathrm{red}})\).  

One can associate to \((P, \mathcal{H}, g)\) a canonical curvature operator \(\mathcal{R}^{P}\) in the sense of Agrachev–Zelenko, defined along normal geodesics via Jacobi curves in the sub‑Riemannian cotangent bundle. Let \(\mathcal{R}^{M}\) be the corresponding curvature operator on \((M, \mathcal{H}_{\mathrm{red}}, g_{\mathrm{red}})\).

**Open Problem (Precise version):**  
Under the above setting, conjecture that there exists a canonical relation between the lifted Jacobi curves on \(P\) and their projections on \(M\) such that  
\[
\operatorname{Tr}(\mathcal{R}^{M}_\lambda) = \operatorname{Tr}(\mathcal{R}^{P}_{\tilde{\lambda}}) - \Phi_G(\lambda),
\]
for some explicitly computable nonnegative function \(\Phi_G\) on the reduced cotangent bundle \(T^*M\) depending only on the differential of the \(G\)-action.  

Determine whether such a formula holds (even conjecturally) when \(P\) is Sasakian and \(G\) acts by contact isometries, and whether \(\Phi_G\) can be interpreted as a “curvature defect” of the symmetry reduction in analogy with O’Neill’s tensor in Riemannian submersions.

Why is it a "good" problem:  
This version is sharply posed: it proposes the existence of a specific quantitative relationship between canonical sub‑Riemannian curvature operators before and after symmetry reduction. It connects deep analytic invariants (Jacobi curves, curvature operators) with geometric group actions. The question generalizes the classical Riemannian O’Neill formulas to the nonholonomic setting, where no Levi–Civita connection exists. Establishing (or disproving) such a formula would provide a fundamental bridge between sub‑Riemannian geometry and symplectic/variational reduction theory. Its apparent simplicity (a trace identity) hides a highly nontrivial interaction between horizontal geometry, curvature, and symmetry. This blend of conceptual clarity, cross‑disciplinary integration, and genuine mathematical depth makes it a “good” problem.